\title{Trending topics using the Twitter-API}
\author{
        Marco Pock-Steen Fraile \\
        Christian Palmh\o j Nielsen \\
        IT University of Copenhagen\\
}
\date{\today}

\documentclass[12pt]{article}

\usepackage{graphicx}
\usepackage{algpseudocode}
\usepackage{algorithm}
\usepackage{amsmath}
\usepackage{listings}


\lstset{
  language=Python,
  showstringspaces=false,
  formfeed=\newpage,
  tabsize=2,
  commentstyle=\itshape,
  basicstyle=\scriptsize,
  morekeywords={models, lambda, forms}
}

\newcommand{\code}[2]{
  \hrulefill
  \subsection*{#1}
  \lstinputlisting{#2}
  \vspace{2em}
}


\begin{document}
\maketitle

\begin{abstract}
Social networks such as Twitter produce enourmeous amounts user-generated data that is a great source for data-mining algorithms to extract information about rising trends and opinions. We have implemented the Misra-Gries algorithm as a way of identifying trending topics in twitter. With the K-minimum values algorithm we can determine if a trending topics popularity is caused by many distinct mentions or is a result of a high number of retweets. We explain how we adapt the Twitter data stream to accomodate the mentioned algorithms.
\end{abstract}

\section{Introduction}
\label{intro}
We wish to make use of the Twitter streaming API to implement our own version of trending topics using the Misra-Gries algorithm, and for the trending topics we would track how many unique tweets have been made per topic, using one of the algorithms for counting distinct elements. We could also count number of tweets for trending topics to compare actual tweets with unique tweets, but this is not algorithmically a challenge.
\newline\newline
Challenges would include finding a suitable k for the Misra-Gries that is a good trade-off between memory space and making sure we take into account the tweets inbetween the actual trending topics that might ``flood'' them out of the array.
\newline\newline
For the counting algorithm we similarly need to find a good hashing algorithm for tweets, and based on this find a good number for kth-minimum distances that gives us a realistic count for unique tweets.

\paragraph{Outline}
The remainder of this article is organized as follows.
Section~\ref{twitter} gives an introduction to the Twitter terminology and ways users can interact with eachother. The streaming model is described in Section~\ref{related work}, followed by the algorithm design and analysis in Section~\ref{algo-data}. Our results are described in Section~\ref{results}. Finally, Section~\ref{conclusions} gives the conclusions.

\section{Twitter}
\label{twitter}
Twitter is a real-time information network that connects you to the latest stories, ideas, opinions and news about what you find interesting. All information on Twitter is distributed through ``tweets'', small 140-character messages where a person can shortly state an opinion, an event or random thoughts. Because of the format, twitter encourages frequent updates from users and as of October 26th, 2012 around 500 million tweets are written per day (24 hours)\cite{Cnet1}. With 140 million users, this means each user, on average, posts 3.5 tweets per day. A typical tweet can be seen below.
\\\\
\includegraphics[width=125mm]{tweet.png}
\\\\
Each user on twitter chooses other users to follow and see an aggregated list of all their tweets. Tweets can be replied to by in the form of ``@name message'', as examplified below. For our purpose, we do not connect replies with previous tweets, but Twitter does provide this in the tweet data.

\begin{quote}
    \emph{@pocky1000 You are totally right. \#ITU is super cool.}
\end{quote}

Users of twitter can also retweet a tweet in their stream. On retweet the user basically shares the same message to their own followers, giving credit to the creator of the message. However, the content of the tweet is identical to the original tweet. An example of a re-tweet follows.
\\\\
\includegraphics[width=125mm]{retweet.png}
\\\\
\subsection{Hashtags and topics}
For the purpose of categorizing tweets, users put hashtags within their tweets. Examples include, as seen above, \#ITU, \#trending and \#Misra. For the purpose of this paper, we will consider these hashtags as topics, and will only consider tweets that contain at least one hashtag. If a tweet contains several hashtags, it will add to our counter for each of them seperately. It will also be used individually for each one when considering uniqueness (so it might be unique to one hashtag but not to another).

A trending topic is defined by a topic that has a high frequency, in other words many people are ``tweeting'' about the same topic at some point in time. A trending topic can either be many unique tweets about the same topic, e.g. if people are tweeting to promote some beneficial cause, or it can be the same message retweeted by many users. We use the public Twitter API to create a stream of which corresponds to a sample of all tweets. 

\subsection{Retrieving the data}
To retrieve the data, we make use of Twitter's streaming API\footnote{https://dev.twitter.com/docs/streaming-apis/streams/public}. They offer two endpoints for getting a live feed of data: ``firehose'' and ``sample''. Firehose feeds all tweets made to twitter, but is not available for general purpose usage. As such we will use the sample, which in our results gives around 500.000 tweets (with hashtags) per day, roughly 1/1000th of the total stream. See more in section \ref{frequent-determine}.


\section{Streaming algorithms}\label{related work}
\subsection{Introduction}\label{algo-intro}
\begin{quote}\textit{
Twitter and other social networks are structured to accomodate personal communication across large networks of friends, and a such produce enormous amounts of data. The open availability of this data through developer APIs makes it's an interesting source for useful real-time information extraction using streaming algorithms \cite{genderprediction}. 
}
\end{quote}

This section introduces the concepts of algorithms that can compute some function of a massively long input stream $\sigma$ such as all public available tweets. In our model this is formalized as a sequence $\sigma = \langle a_{1}, a_{2},...,a_{m}\rangle $, where the elements of the sequence (in our case $tweets$) are drawn from the universe $[n] =: \{1, 2,..., n\}$. Note the two size parameters: the stream length, $m$, and the universe size, $n$.

\subsection{Goal}\label{algo-goals}
Our goal is to find the trending topics on twitter over a given time period. The idea is that our algorithm will be optimized for different intervals, eg. 1-hour and 24-hour periods, and that at any given time it can be queried to get the actual trending topics for previous interval (assuming the stream has been running for that minimum amount of time).
Our goal will be to process the input stream using a small amount of space $s$, i.e., to use s bits of random-access working memory. Since $m$ and $n$ are to be thought of ``huge'' we want to make $s$ much smaller than these. Ideally we want to achieve $s$ = O(log $m$ + log $n$), because this is the amount of space needed to store a constant number of elements from the stream and a constant number of counters that can count up to the length of the stream. \cite{Amit}

Specifically we want to analyze the Twitter input stream to find tags that that have a high occurence over some period of time, and be able to return these e.g. as a top-10 list for both 1 hour and 24 hour periods. Besides returning the most popular topics, we wish to be able to say something about if the topic has a high frequency because many different people have mentioned it, or it has been retweeted a lot of times.

\subsection{Finding frequent items}\label{algo-frequent}
Ideally, we want to find the frequency of topics in $tweets$ in our input stream $\sigma = \langle a_{1}, a_{2},...,a_{m}\rangle$. We define a frequency vector $f = (f_{1},...,f_{n})$, where $f_{1} + ... + f_{n} = m$. 


To find the exact frequency of each topic would require that we store the exact number of times it has been seen, and with the vast number of topics seen (see section \ref{frequent-determine}) it would require a large amount of space to store. Instead, we can find the most frequent items seen over a period of time using the \textit{Misra-Gries} algorithm.\cite{Amit}
\newline

From the algorithm description we have $f_j - \frac{m}{k} \leq \hat{f_j} \leq f_j$. From this we can derive that any item with a $\hat{f_j} > \frac{m}{k}$ will be in our map $A[j]$ with a positive score, and we simply need to make sure the topics we wish to capture for a given time period are above the $\frac{m}{k}$ threshold. We will discover $m$ and derive $k$ from it in the following section.

\subsubsection{Determine k and m for Misra-Gries}
\label{frequent-determine}
In order for us to determine a suitable $k$ we will first have to find $m$ for a given timeframe. We wrote a piece of code that for 1 hour and 24 hours would count accurately \textit{a)} The total number of tweets, \textit{b)} The number of distinct topics and \textit{c)} The number of occurrances for each topic. We ran this a few times on different days to get an average on the number of tweet and occurances of the top 10 topics. An example of each sampling (1 hour and 24 hour respectively) can be seen below.\newline

Total number of tweets in 60 minutes: 26733, with 19117 unique topics\\

$\begin{array}{ll}
    454: & \#PeoplesChoice \\
    446: & \#StarAc \\
    331: & \#IHaveACrushOn \\
    316: & \#tbt \\
    265: & \#gameinsight \\
    251: & \#RT \\
    240: & \#TeamFollowBack \\
    223: & \#LouisAndBoris \\
    146: & \#musicfans \\
    140: & \#FF \\
\end{array}$
\\
\\

Total number of tweets in 1440 minutes: 534408, with 214345 unique topics\\

$\begin{array}{ll}
    17064: & \#FF \\
    7271: & \#gameinsight \\
    6589: & \#PeoplesChoice \\
    4775: & \#TeamFollowBack \\
    4433: & \#RT \\
    3692: & \#androidgames \\
    3522: & \#TFBJP \\
    3510: & \#MessageToMyEx \\
    3456: & \#ff \\
    3370: & \#android \\
\end{array}$
\\
\\
Given this data, we can determine an approximation for $k$. We want to include at least the top 10 topics for the time periods, so we maintain two seperate data structures for each. For the 60 minute period, we want topics with a minimum frequency of \textbf{140} occurrences over \textbf{26733} total events. This means we have $k = \frac{26733}{140} \approx 190$. If we take the average of all our samplings, we find the average to be $k = 155$ for a 24-hour period, and $k = 165$ for a 1-hour period. We assumed that $k$ would be larger for the 24 hour period, but as some topics keep repeating frequently spread over a day, they actually require a smaller k for bigger intervals. Also, our samples are not for longer periods (months), so this might not be completely accurate.
\\
\\
We discuss the results and our conclusions in section \ref{trending-discussion}


\subsection{Finding distinct values}\label{algo-distinct}
K-minimum values (KMV) is is a probalistic distinct value counting algorithm, that is intuitive and easy to implement \cite{kmv}. Suppose we have a good hash function that return evenly distributed values in the hash space $[0-1]$, then you could estimate the number of distinct values you have seen by knowing the average spacing between values in the hash space. The main challenge is to find a good hash function, and to select the number of minimum $k$ values on which to approximate the average spacing. If the hash values were indeed evenly distributed, we could keep just keep track of the minimum value, a get a good estimate of distinct values. However taking only one value opens up to a lot of variance and would rely heavily upon the``goodness'' of the hash function. In order to improve this Bar-Yossef\cite{Bar-Yossef} suggests keeping the k-smallest values, to give a more realistic estimate.
   
Other examples of the usages of data stream algorithms are described in this chapter.
\subsection{Implementation of algorithms}

\subsubsection{Misra-Gries Algorithm}\label{misra-gries}
The algorithm first initializes a dictionary with $k$ number of values. The keys in the dictionary are elements seen in the stream, and the value are counters associated with the elements. Then there is a process function that is executed each time we see a new element. If a new element is already in the dictionary, it's value will be increased by 1, otherwise if the number of elements in $A$ is less than $k$, the element will be inserted an its value set to 1. If the length of $A$ is equal to $k$, all values are decreased by 1, and removed if the value is equal to 0 \cite{Amit}. Finally we return the key value pairs with the highest frequencies.

We use Misra-Gries to find frequent $topics$ in our data stream of tweets using a one-pass algorithm, see pseudo-code \ref{misra-pseudo} for details.
\begin{algorithm}
\caption{Misra-Gries Algorithm}\label{misra-pseudo}
\begin{algorithmic}[1]
\State $A\gets Initialize Array$

\Function{Process}{$j$}
\If {$j \in keys(A)$}
    \State $A[j] \leftarrow A[j] + 1$
\ElsIf {$|keys(A)| < k-1 $}
    \State $A[j] \leftarrow 1$
\Else
    \For{$l \in keys(A)$}
        \State {$A[l] \leftarrow A[l] - 1$}
        \If {$A[l] = 0$} remove $l$ from A
        \EndIf
    \EndFor
\EndIf
\EndFunction
\end{algorithmic}
\end{algorithm}
\subsubsection{Data structure}
When a new tweet is seen the algorithm has to determine if the contained topic has been seen before. We use a dictionary with key-value pair as topic-frequency. The dictionary provides $O(1)$ time lookup and increment operations. When we decrement all elements by one we iterate over the dictionary which takes $O(k)$ time. Finally we sort the dictionary, but since it's limited to $k$ items, this is not a very time consuming operation.
\subsubsection{Data set}\label{algo-data}

\subsubsection{Analysis of algorithm}\label{algo-analysis}
Misra-Gries uses one pass with worst case running time $O(k)$ for each element of the stream $[n]$. Insertion and increments of a topic takes $O(1)$ time if the length of the dictionary is less than $k$. However the worst case scenario would decrement all values of the dictionary every time the dictionary is full, a maximum of $\frac{k}{n}$ times, so we can argue that the amortized running time for each process is $O(1)$.


\section{Results}\label{results}
In this section we describe the results of running the algorithm for a 1 hour period using the found values of \textit{k} for \textit{Misra-Gries} and \textit{KMV} respectively.
\\\\
$\begin{array}{ll}
    902: & \#peopleschoice (distinct: 922.0) \\
    575: & \#breakoutartist (distinct: 637.0) \\
    326: & \#cumannanyayap (distinct: 343.0) \\
    180: & \#teamfollowback (distinct: 236.0) \\
    160: & \#onedirection (distinct: 231.0) \\
    125: & \#camilafollowspree (distinct: 171.0) \\
    72: & \#gameinsight (distinct: 227.0) \\
    35: & \#tfbjp (distinct: 106.0) \\
    24: & \#android (distinct: 139.0) \\
    20: & \#androidgames (distinct: 137.0) \\
\end{array}$
\\\\
This result shows why it might be relevant to consider distinct tweets when finding trending topics. Despite being mentioned more frequently, \#bfbjp and \#camillafollowspree both have fewer distinct tweets than the respective topics below them.

\subsection{Comparison with Twitters trending topics}\label{twitter-result}
The results may differ from Twitters official trending topics for various reasons, which include but are not limited to the following: 

\begin{enumerate}
	\item Twitters trend emphasizes novelty over popularity. This means that topics that have a high increase in frenquency over a short period of time are rated higher than topics that are always popular like \#android and \#gameinsight. 
	\item We only have access to a sample of all public tweets which is not guaranteed to be statistical accurate representation.
    \item For simplicity of implementation we only consider tweets that contains one or more tags. This does not seem to be the case for twitters own trending topics.
\end{enumerate}

By a quick comparison with the top-10 trending topics on Twitter with our results, it shows that there indeed exists some overlap of topics, but it is hard to define how much precisely.




\section{Discussion}\label{conclusions}

\subsection{Trending vs popular}\label{trending-discussion}
From our results, we can see that between our 1-hour and 24-hour data mining period a lot of the same topics were in the top10. One thing $Misra-Gries$ does not take into account is whether a topic is new or old within a given period - if a topic occurs more frequently than our threshold at the beginning of our time period, it will remain there for the remainder of the period. An alternative to $Misra-Gries$ would be the $Lossy Counting$ algorithm. It works in similar ways, but topics found earlier in the stream are more likely to get pruned than topics found later in the steam, which gives topics is mentioned over a very short period of time a higher value than topics that have been less, but more regularly, mentioned over a longer period.

We can conclude that our chosen algorithm 

\subsection{Effect of distinct tweets}\label{trending-distinct}

We worked hard, and achieved very little.

\section{Appendix: Code}

\code{Base code}{base.py}

\code{Data Sampling}{sampling.py}

\code{Trending topics}{trends.py}

\bibliographystyle{abbrv}
\bibliography{simple}

\end{document}
This is never printed